\section{Country-Level Findings}

\subsection{Demand scale concentrates in a small number of markets}

The platform-tourism panel is dominated by large Western and Mediterranean markets. France, Spain, and Italy are the key scale markets. Even modest shoulder-season shifts in these countries may have material commercial value because the absolute demand base is large.

\begin{finding}
Scale is a necessary first filter. Markets with very large demand bases can justify deeper regional analysis even when their national seasonality pattern is less extreme than smaller Mediterranean markets.
\end{finding}

\reportfigure{top_demand.pdf}{0.95\textwidth}{Top countries by recent platform-tourism demand. The business interpretation is that scale-led markets require regional disaggregation before campaign conclusions can be drawn.}

\subsection{Mediterranean markets show the clearest summer concentration}

The seasonality heatmap normalizes each country to its own annual total, so it shows timing rather than absolute size. Greece and Croatia show the clearest summer concentration pattern, while larger markets such as France, Italy, and Spain require additional decomposition because national averages likely hide strong regional differences.

\begin{finding}
The strongest theoretical shoulder-season signal comes from countries where summer demand is concentrated but shoulder months are not at winter-trough levels. This suggests that demand exists outside the peak but may need destination-specific activation.
\end{finding}

\reportfigure{seasonality_heatmap.pdf}{0.95\textwidth}{Within-year concentration of platform nights. Rows sum to 100 percent, so the figure shows seasonal shape rather than market size.}

\subsection{Growth separates watchlist markets from investable scale markets}

High percentage growth in small markets should be interpreted carefully. A smaller country can show rapid growth from a low base without yet justifying the same follow-up effort as a large or highly seasonal market.

\begin{finding}
Growth is useful as a differentiating signal but should not dominate the business decision. High-growth small markets are better treated as watchlist candidates unless they also show sufficient demand scale or clear seasonality opportunity.
\end{finding}

\reportfigure{growth_vs_demand.pdf}{0.88\textwidth}{Growth versus demand level. The useful business distinction is between scale-led candidates, seasonality-led candidates, and small high-growth watchlist markets.}

\subsection{The opportunity ranking produces a focused shortlist}

The default scoring model prioritizes five countries for Phase 2: \countrycode{EL}, \countrycode{FR}, \countrycode{IT}, \countrycode{HR}, and \countrycode{ES}. The order is less important than the shortlist itself because each country ranks highly for a different reason.

\reportfigure{opportunity_ranking.pdf}{0.88\textwidth}{Top countries by composite opportunity score. The score ranks countries for follow-up investigation rather than direct budget allocation.}

\reportfigure{opportunity_components.pdf}{0.88\textwidth}{Component decomposition for top-ranked countries. This explains whether a country ranks because of scale, growth, seasonality, or a balanced profile.}

\begin{finding}
The top-five markets are not homogeneous. Greece and Croatia are seasonality-led; France and Spain are scale-led; Italy is mixed. Phase 2 should therefore adapt the analysis brief by country instead of treating the shortlist as one uniform segment.
\end{finding}
